\documentclass[12pt, letterpaper]{article}

% --- Paquetes básicos ---
\usepackage[spanish, es-tabla]{babel}
\usepackage[utf8]{inputenc}
\usepackage[T1]{fontenc}
\usepackage{newpxtext,newpxmath} % Fuente estilo Palatino
\usepackage[margin=3cm]{geometry}
\usepackage{titlesec}
\usepackage[autostyle]{csquotes}

% --- Para dibujar la V de Gowin ---
\usepackage{tikz}
\usetikzlibrary{positioning,babel,calc}

% --- Configuración de la sección (1. IDEA DE...) ---
\titleformat{\section}
  {\normalfont\large\scshape\centering} 
  {\thesection.}
  {0.5em}
  {}

\begin{document}

% --- Encabezado Manual ---
\begin{center}
    \vspace{0.8cm}
    {\Large \bfseries DÍGITO 1 LIDERA FRECUENCIA PERO BENFORD PREDICE EL DOBLE DE SU PRESENCIA} \\
    \vspace{0.4cm}
\end{center}

\vspace{1.5cm}

\noindent
\textbf{Nivel descriptivo} \newline

\noindent
Titular: Dígito 1 lidera frecuencia pero Benford predice el doble de su presencia. \newline

\noindent
Nombre del hallazgo: Desviación sistemática de la distribución de primeros dígitos respecto a la Ley de Benford. \newline

\noindent
Resumen en una oración: Los primeros dígitos en contraseñas humanas siguen una tendencia similar a Benford pero significativamente más aplanada. \newline

\noindent
Método o análisis que lo produjo: Comparación directa entre la frecuencia observada del primer dígito en \texttt{rockyoubenford} y la distribución teórica de la Ley de Benford, visualizada en el Gráfico 1. \newline

\noindent
Evidencia: El Gráfico 1 (Figura 2.6) muestra que el dígito 1 aparece como primer dígito en aproximadamente el 18\% de los casos, mientras que la Ley de Benford predice cerca del 30\%. Los dígitos del 3 al 9 aparecen con frecuencias más parejas entre sí de lo que la ley esperaría, y el dígito 0, que Benford no contempla como primer dígito significativo, aparece con una frecuencia destacada de aproximadamente el 28\% según el Cuadro 1 (Figura 2.3).
\newpage

\noindent
\textbf{Nivel analítico} \newline

\noindent
Conexión con la pregunta de investigación: La desviación respecto a Benford confirma que las contraseñas humanas no siguen ni la aleatoriedad uniforme ni los patrones matemáticos naturales, sino una lógica cultural propia basada en fechas, años y secuencias conocidas. Esta lógica cultural es precisamente lo que un bot no replicaría, lo que la convierte en una variable discriminante con fundamento teórico sólido. \newline

\noindent
Contraste con la literatura: Este resultado es coherente con lo planteado por Caputi (2019) en \textit{Ley de Benford}, quien describe que la ley se cumple en datos de origen natural pero puede desviarse cuando los datos responden a convenciones culturales. También conecta con Olsen (2018), quien aplicó la Ley de Benford al primer dígito en su análisis de contraseñas y encontró patrones similares de desviación. \newline

\noindent
Lo que NO explica este resultado: Este hallazgo no determina si la desviación es estadísticamente significativa mediante una prueba formal como chi-cuadrado o MAD, ni si el patrón se mantiene al segmentar por longitud de contraseña o por chunk. Tampoco explica por qué el dígito 0 tiene una presencia tan elevada como primer carácter, cuando Benford lo excluye por definición. \newline

\noindent
Implicación para el siguiente paso: En la Bitácora \#3 se deberá aplicar una prueba de bondad de ajuste formal (chi-cuadrado o MAD siguiendo la metodología de Caputi) para cuantificar el grado de desviación respecto a Benford y determinar si esta diferencia es suficientemente grande como para usarse como criterio discriminante entre contraseñas humanas y generadas por bots.

\end{document}
